\documentclass[11pt]{article}
\usepackage{amsmath,amsthm,amssymb}
\newcommand{\R}{\mathbb{R}}
\newcommand{\Lap}[1]{\Delta_{#1}}
\DeclareMathOperator{\Deg}{deg}
\theoremstyle{plain}
\newtheorem{theorem}{Theorem}[section]
\newtheorem{lemma}[theorem]{Lemma}
\newtheorem{proposition}[theorem]{Proposition}
\newtheorem{corollary}[theorem]{Corollary}
\theoremstyle{definition}
\newtheorem{definition}[theorem]{Definition}
\theoremstyle{remark}
\newtheorem*{remark}{Remark}
\title{Discrete Harmonic Functions on Finite Weighted Graphs}
\begin{document}
\maketitle
\begin{abstract}
We establish a maximum principle and an energy identity for discrete harmonic functions on finite connected weighted graphs.
\end{abstract}
\section{Graph Laplacian and harmonicity}\label{sec:setup}
Let $G=(V,E)$ be a finite connected graph with positive edge weights $w_{uv}=w_{vu}>0$.
\begin{definition}[Weighted Laplacian]\label{def:lap}
For any function $f:V\to\R$, the weighted Laplacian at $u\in V$ is
\begin{equation}\label{eq:lap}
(\Lap{G} f)(u) = \sum_{v\sim u} w_{uv}(f(u) - f(v)).
\end{equation}
We call $f$ harmonic on $U\subseteq V$ if $\Lap{G} f(u)=0$ for every $u\in U$.
\end{definition}
\begin{lemma}[Dirichlet form identity \cite[Proposition 2.1]{keller2020}]\label{lem:dirichlet}
For all $f,g:V\to\R$,
\[
\sum_{u\in V} f(u)(\Lap{G} g)(u) = \frac{1}{2}\sum_{u\sim v} w_{uv}(f(u)-f(v))(g(u)-g(v)).
\]
\end{lemma}
\begin{proof}
Reindexing the sum by swapping $u$ and $v$ and using symmetry $w_{uv}=w_{vu}$ yields the identity directly.
\end{proof}
\section{Maximum principle and uniqueness}\label{sec:max}
\begin{theorem}[Discrete maximum principle]\label{thm:max}
Let $U\subsetneq V$ be a nonempty proper subset. If $f:V\to\R$ is harmonic on $U$, then
\begin{equation}\label{eq:max-bound}
\max_{u\in U} f(u) \le \max_{v\in V\setminus U} f(v).
\end{equation}
\end{theorem}
\begin{proof}
Suppose the maximum is attained at an interior vertex $u_0\in U$. Since $f(u_0)$ is the weighted average of its neighbours, every neighbour attains the same value. Connectivity propagates this value to $V\setminus U$.
\end{proof}
\begin{corollary}[Uniqueness of the Dirichlet problem]\label{cor:unique}
Given boundary values on $V\setminus U$, there is at most one function harmonic on $U$ agreeing with those values.
\end{corollary}
\begin{remark}
The positivity hypothesis $w_{uv}>0$ cannot be relaxed to nonnegative weights without risking disconnected components.
\end{remark}
\end{document}
